\scanpage{6}\mspage{4}

$u_{2n-i}$ est alors la \deleted{\unsure{ques|sous}} conjonction des suivantes :
\[
\left[\begin{aligned}
 (3.1.)_i\quad &(u\in \unsure{Z|\mathfrak Z|\mathfrak z}\unsure{E|\mathcal E}^0,\ \deleted{\unsure{u\pi_\alpha=0\ \text{pour }\alpha<i},}\ 
 \text{\textcircled{$E_i^{2n-i}$}}u=0)\Longrightarrow \deleted{\unsure{u=0}}\inserted{\pi_i u=0},\\
 (3.2)_i\quad &(u\in ZE^0,\ \deleted{\unsure{\pi_\alpha u=0\ \text{pour }\alpha\ne 2n-i},}\ uE_i^{2n-i}=0)
 \Longrightarrow\deleted{\unsure{u=0}}\inserted{\pi_{2n-i}u=0}
\end{aligned}\right]
\]
En effet, la relation (3.1.) implique que \deleted{\unsure{$\pi_i(ZE^0)$|$\pi_i\in ZE^0$}} \inserted{l'annulateur de $\pi_i$ dans $ZE^0$} est l'idéal des $u\in ZE^0$ tels que $E_i^{2n-i}u=0$ ; or dans un anneau commutatif, un idempotent est connu quand on connaît son annulateur. De même pour $\pi_{2n-i}$, son annulateur dans $ZE^0$ est formé des $u\in ZE^0$ tels que $uE_i^{2n-i}=0$.

En même temps, les conditions $(3.1)_i$, $(3.2.)_i$ pour \unsure{tt|cet} $i$, donnent un \deleted{condition} procédé pour déterminer de proche en proche les $\pi_i$.

\medskip\hrule\medskip

4. Lorsque $E^*$ est comme dans 1., une condition suffisante pour avoir $(3.1)_i$ \inserted{resp. $(3.2)_i$} est la \deleted{\unsure{suff}} suivante

$(4.1)_i$ L'intersection des noyaux des $\unsure{v|u}:M^i\longrightarrow M^{2n-i}$, pour $v\in E_i^{2n-i}$, est réduite à 0.

$(4.2)_i$ La somme des images des $v:M^i\unsure{-|\to}M^{2n-i}$, pour $v\in E_i^{2n-i}$, est \deleted{\unsure{se}} égale à $M^{2n-i}$.

Ces deux conditions sont satisfaites si on a

$(4.3)_i$ Il existe \inserted{$v:M^i-M^{2n-i}$,} $v\in E_i^{2n-i}$, qui soit un \unsure{isom$^{\mathrm e}$|isom.} de
% archive sheet number: 5
